\section{Introduction}
Modern industrial infrastructures rely heavily on rotating machinery, such as wind turbines, aero-engines, and high-speed railways, where critical components including bearings and gearboxes operate under demanding service conditions \cite{zhai2021multi}. The health status of these machines is critical to operational safety and economic efficiency, as unforeseen failures often lead to catastrophic accidents and large financial losses \cite{lei2020applications}. Consequently, Intelligent Fault Diagnosis (IFD) has emerged as a cornerstone of smart manufacturing, using sensor-based condition monitoring to detect anomalies before they escalate into system-wide failures. However, as the complexity of industrial systems increases, the demand for more robust, autonomous, and high-precision diagnostic frameworks is intensifying to ensure the continuity of critical production lines \cite{li2020deep}. Fig.~\ref{fig:concept} illustrates the conceptual overview of the proposed cross-modality semantic alignment approach for industrial zero-shot fault diagnosis.

Despite the advancements of deep learning in fault diagnosis, the majority of existing methods operate under a closed-set assumption, which requires intensive manual labeling and a balanced distribution of training samples across all predefined health states \cite{zhang2022digital}. In practical industrial scenarios, machinery typically operates in normal states for the vast majority of its lifespan; therefore, fault samples---particularly those representing specific early-stage degradation---are rare and difficult to acquire \cite{liu2018zero}. This imbalance creates a significant long-tail distribution problem where traditional supervised models fail to generalize to unseen fault categories. To address this data scarcity and the impracticality of exhaustive labeling, research focus has shifted from standard supervised learning toward Zero-Shot Learning (ZSL). ZSL aims to recognize new fault classes without physical training samples by using auxiliary information to bridge the gap between known and unknown categories \cite{wang2019zero, socher2013zero}.

Current ZSL approaches in fault diagnosis primarily rely on semantic embeddings, such as manual attributes or word vectors (e.g., Word2Vec or GloVe), to associate vibration signals with class labels \cite{lampert2014attribute}. However, these methods face three critical limitations. First, the definition of manual attributes is labor-intensive and inherently subjective, requiring domain experts to enumerate physical characteristics for every possible fault type \cite{fu2018zero}. Second, generic word embeddings are typically trained on web-scale natural language corpora, such as Wikipedia, which lack the technical specificity required for industrial diagnostics. For instance, the term ``outer race fault'' carries distinct engineering connotations that generic vectors fail to capture with sufficient depth \cite{mikolov2013efficient}. Third, existing models often neglect the complex temporal dependencies inherent in high-frequency vibration data, leading to a domain gap where the alignment between signal features and the abstract semantic space remains brittle \cite{xian2018zero, buchert2020zero}.

To overcome these challenges, this paper proposes the Cross-Modality Semantic Alignment Transformer (CMSA-Trans), a new zero-shot fault diagnosis framework that uses Large Language Models (LLMs) to bridge the semantic divide. Unlike traditional methods, CMSA-Trans uses the descriptive capabilities of LLMs to generate high-fidelity, domain-specific fault semantics, thereby transforming abstract class labels into detailed technical descriptions. By formulating fault diagnosis as a cross-modal alignment task, the framework employs a Transformer-based temporal encoder to extract discriminative features from raw vibration signals. These features are subsequently projected into a shared latent space along with the LLM-generated semantics. This alignment ensures that the model learns the underlying physical relationships between signal patterns and their linguistic definitions, enabling the recognition of unseen faults through the semantic proximity of the signal to known physical phenomena \cite{vaswani2017attention, radford2021learning}.

The primary contributions of this research are summarized as follows. First, we introduce the CMSA-Trans framework, which reformulates zero-shot fault diagnosis as a cross-modal semantic alignment task, effectively bypassing the requirement for labeled samples of target fault categories. Second, we propose a systematic method to use LLMs for the generation of expert-level fault semantics, capturing the nuanced physical characteristics of rotating machinery faults more effectively than traditional word embeddings. Third, a Transformer-based temporal feature extractor is designed to capture long-range dependencies in vibration signals, providing a robust representation for alignment with high-dimensional semantic descriptors. Finally, extensive experiments on multiple open-source and industrial datasets show that CMSA-Trans greatly outperforms state-of-the-art ZSL methods, achieving high accuracy on unseen fault classes and showing remarkable generalization across diverse operating conditions.

Also, the integration of LLMs, such as GPT-4, represents a approach shift in semantic feature engineering for industrial applications \cite{bubeck2023sparks}. Unlike static word embeddings that aggregate broad linguistic associations, LLMs can be prompted to synthesize multi-faceted technical descriptions that capture the underlying physics of a fault, such as impact repetition frequencies and localized high-frequency resonances \cite{brown2020language}. This approach allows for a richer and more discriminative semantic space, where the distance between fault classes reflects actual physical mechanisms rather than simple linguistic co-occurrence in general-purpose text. By harnessing these LLM-generated semantics, CMSA-Trans provides a data-efficient and highly descriptive alternative to traditional ZSL approaches.

The importance of robust cross-modal alignment is critical. In traditional deep learning models, features extracted from sensor data are typically mapped directly to one-hot encoded labels, which contain no inherent physical information and thus cannot generalize to unseen fault types \cite{goodfellow2016deep, lecun2015deep}. In contrast, our semantic alignment approach uses a projection layer to map the high-level signal representations of the Transformer into the same dimensional space as the LLM-derived technical descriptors. This ensures that the model learns to associate specific signal patterns, such as periodic transients or broadband noise increases, with physical descriptions. When a new fault type with a unique technical signature occurs, the model can identify it by matching the signal pattern to the linguistic definition of the fault, even in the absence of labeled training examples.

Recent advancements in attention-based architectures have transformed the processing of sequential data in various fields, ranging from natural language processing to computer vision \cite{devlin2018bert, dosovitskiy2020image}. In the context of rotating machinery, vibration signals are inherently time-variant and exhibit long-range dependencies, where early-stage fault signatures may be obscured by noise or periodic components from normal operation \cite{zhang2021attention}. While Recurrent Neural Networks (RNNs) and Long Short-Term Memory (LSTM) units have been employed to handle such data, they often struggle with gradient vanishing or capturing dependencies over the very long windows typical of high-frequency samples \cite{hochreiter1997long}. The self-attention mechanism of the Transformer provides a more effective method to weight different time steps according to their relevance to the diagnostic task, allowing the CMSA-Trans model to focus on the most informative transients and spectral characteristics across the entire signal window.

The shift toward using large-scale pre-trained models for industrial semantic understanding is motivated by the extensive knowledge encapsulated within these models \cite{zhao2023brain}. While an engineer might define a bearing fault by its physical dimensions, an LLM can provide a full textual profile that links those dimensions to the expected sensor responses and characteristic frequencies. By aligning these rich descriptions with raw data, our approach moves closer to an expert-informed diagnostic system that mimics the reasoning of an experienced technician. This capability is particularly vital for critical infrastructure where failures are rare but catastrophic, as the ability to correctly identify a new failure mode upon its first occurrence is essential for preventing downtime and maintaining structural integrity.

Expanding the scope of zero-shot learning to include cross-modality also addresses the hubness problem often encountered in semantic embedding spaces, where certain vectors become nearest neighbors to many unrelated items in high-dimensional space \cite{radovanovic2010hubs}. By using a more descriptive and technically grounded semantic space provided by LLMs, the separation between distinct fault modes---such as inner-race versus outer-race bearing defects---becomes more pronounced. This improved class separation in the semantic domain translates directly to improved diagnostic reliability and higher confidence scores for unseen categories. Validation on a variety of machinery types, including complex planetary gearsets and varying load-speed conditions, further confirms that the CMSA-Trans framework provides a scalable and robust solution for the next generation of intelligent PHM (Prognostics and Health Management) systems. These advantages collectively position CMSA-Trans as a practical and theoretically grounded alternative to conventional supervised diagnostic paradigms.

The remainder of this paper is organized as follows. Section II provides a full review of the related literature regarding intelligent fault diagnosis and zero-shot learning \cite{li2022zero}. Section III details the proposed CMSA-Trans architecture, including the process of LLM semantic generation and the cross-modal alignment mechanism. Section IV presents the experimental setup, results, and comparative analysis. Finally, Section V concludes the paper and discusses future research directions.
